
\section{Task Setup and Datasets} %




\label{sec:task_def}

Our task is formalized as follows:
Given a query $q$ and a corpus of text passages $\mathcal{D}$,
the system %
is required to return an output $\mathcal{S}$%
, which consists of $n$ statements $s_1$, ..., $s_n$, 
and each statement $s_i$ cites a list of passages $\mathcal{C}_i = \{c_{i,1}, c_{i, 2}, \ldots\}$\footnote{In practice, we allow at most 3 citations for each statement as more citations usually do not help.}, where $c_{i,j}\in \mathcal{D}$. 
In this work,
we segment LLMs' output into statements by sentence boundaries.\footnote{
    QAMPARI requires a list as the answer, and we choose each entity in the generated list as a statement.} 
While LLMs may include sentences that do not require a citation, such as ``\ti{I'm happy to help}'', 
we observe that almost all sentences that LLMs output  provide valuable information and  require citations, similar to findings in
\citet{liu2023evaluating}. 
In this work, citations are enclosed by box brackets such as \texttt{[1][2]}.

We divide the corpus $\mathcal{D}$ into 100-word passages following previous works on open-domain question answering~\cite{karpukhin-etal-2020-dense,petroni-etal-2021-kilt,piktus2021web},
in contrast to
commercial systems like Bing Chat, which cite entire Web pages. %
We take  100-word passages because
it is easier for humans to verify, and %
allows for more retrieved passages to fit in LLMs' limited context.  %


We choose QA datasets so that 
(1) they contain factual questions, in which references are important; 
(2) questions require long-text answers that cover multiple aspects; %
(3) answering the questions  requires synthesizing multiple sources.
We select three datasets %
(Table~\ref{tab:datasets}) and introduce them below. %
See \S\ref{app_sec:dataset_stats} for additional statistics. 

 
\textbf{ASQA}~\cite{stelmakh2022asqa} is a long-form factoid dataset.
As shown in Figure~\ref{fig:main}, 
each question is an ambiguous question from AmbigQA~\cite{min-etal-2020-ambigqa} that requires multiple short answers to cover different aspects, 
and the dataset provides a long-form answer that covers all short answers.
Since most questions can be answered by Wikipedia, we use the 2018-12-20 Wikipedia snapshot as $\mathcal{D}$.  


\textbf{QAMPARI} \cite{rubin2022qampari}  
is a factoid QA dataset constructed from Wikipedia,
where the answer is a list of entities that are  drawn from different passages.
Same as ASQA, we use the 2018-12-20 Wikipedia as the corpus.%


\textbf{ELI5} \cite{fan-etal-2019-eli5} 
is a long-form QA dataset built on the Reddit forum ``Explain Like I'm Five''.\footnote{\url{https://www.reddit.com/r/explainlikeimfive/}} %
Most ELI5 questions are how/why/what questions that require long answers and multiple passages as evidence.
Due to the diverse topics discussed in the questions, we use
Sphere~\cite{piktus2021web}---a filtered version of Common Crawl\footnote{\url{https://commoncrawl.org}. We also filter out any Web pages from Reddit. }---as the corpus. %
The ELI5 dataset is widely used in related work due to its challenging nature~\cite{nakano2021webgpt,menick2022teaching,liu2023evaluating}.



We randomly select 1,000 examples from the development set of each dataset for \citeeval{}.
Our benchmark primarily assesses the citation capabilities of existing LLMs and does not provide training data,
as
there are no available examples that provide supervision for citations
in these datasets.
